\documentclass[conference]{IEEEtran}
\IEEEoverridecommandlockouts
\usepackage{cite}
\usepackage{amsmath,amssymb,amsfonts}
\usepackage{graphicx}
\usepackage{textcomp}
\usepackage{xcolor}
\usepackage{booktabs}
\usepackage{multirow}
\usepackage{array}
\usepackage{url}

\def\BibTeX{{\rm B\kern-.05em{\sc i\kern-.025em b}\kern-.08em
    T\kern-.1667em\lower.7ex\hbox{E}\kern-.125emX}}

\newcommand{\optionalfigure}[4]{%
\IfFileExists{#1}{%
\begin{figure}[ht]
\centering
\includegraphics[width=\linewidth]{#1}
\caption{#2}
\label{#3}
\end{figure}
}{\typeout{Optional figure #1 not found; skipping.}}%
}

\begin{document}

\title{A Data Preprocessing, Visualization, and Hypothesis-Testing Framework for IPL Performance Analytics}

\author{
\IEEEauthorblockN{Varsha Dange}
\IEEEauthorblockA{\textit{Department of Artificial Intelligence and Machine Learning} \\
\textit{Vishwakarma Institute of Technology} \\
Pune, India \\
varsha.dange1@vit.edu}
\and
\IEEEauthorblockN{Tanmay S. Kumbhare}
\IEEEauthorblockA{\textit{Department of Artificial Intelligence and Machine Learning} \\
\textit{Vishwakarma Institute of Technology} \\
Pune, India \\
tanmay.kumbhare24@vit.edu}
\and
\IEEEauthorblockN{Kasturi Rasekar}
\IEEEauthorblockA{\textit{Department of Artificial Intelligence and Machine Learning} \\
\textit{Vishwakarma Institute of Technology} \\
Pune, India \\
rasekar.kasturi24@vit.edu}
\and
\IEEEauthorblockN{Hemantkumar S. Lakhane}
\IEEEauthorblockA{\textit{Department of Artificial Intelligence and Machine Learning} \\
\textit{Vishwakarma Institute of Technology} \\
Pune, India \\
hemantkumar.lakhane24@vit.edu}
\and
\IEEEauthorblockN{Yash Lakhe}
\IEEEauthorblockA{\textit{Department of Artificial Intelligence and Machine Learning} \\
\textit{Vishwakarma Institute of Technology} \\
Pune, India \\
yash.lakhe24@vit.edu}
\and
\IEEEauthorblockN{Dhanashri Khode}
\IEEEauthorblockA{\textit{Department of Artificial Intelligence and Machine Learning} \\
\textit{Vishwakarma Institute of Technology} \\
Pune, India \\
dhanashri.khode24@vit.edu}
}

\maketitle

\begin{abstract}
The Indian Premier League (IPL) generates high-volume ball-by-ball data, but many student and practitioner analyses remain dashboard-centric and underemphasize data reliability, inferential validation, and domain-correct feature engineering. This paper presents a unified framework that combines preprocessing, descriptive visualization, and hypothesis testing for IPL performance analytics using match-level and delivery-level data from 2008--2019. The proposed pipeline cleans 179,078 raw deliveries and 756 raw matches, removes duplicates, excludes no-result matches and super-over records, normalizes franchise names, and derives phase-aware batting and bowling indicators. The cleaned corpus improves overall data-quality score from 95.44\% to 96.17\%, while reducing franchise-name variants from 15 to 12. The analytical stage combines visual summaries with non-parametric statistical tests. Results show a significant upward trend in season-wise scoring for both first innings ($\rho=0.748$, $p=0.005$) and second innings ($\rho=0.748$, $p=0.005$). Phase-wise run rates differ strongly ($H=586.64$, $p<0.001$), with death overs producing the highest median scoring rate (9.60 runs per over). Team batting strike rate is substantially higher in wins than in losses (137.64 vs. 119.65, $p<0.001$), while winning bowling units exhibit lower economy (7.36 vs. 8.45, $p<0.001$) and higher wicket counts (7.13 vs. 4.50, $p<0.001$). The study demonstrates that statistically grounded cricket analytics requires not only visualization, but also defensible preprocessing and rigorous inferential testing.
\end{abstract}

\begin{IEEEkeywords}
IPL analytics, cricket analytics, data preprocessing, hypothesis testing, sports analytics, data visualization, inferential statistics.
\end{IEEEkeywords}

\section{Introduction}
Modern analytics operates in a regime of data abundance, where the main challenge is no longer data access but reliable transformation into trustworthy evidence. This challenge is particularly visible in professional sports, where event streams, match logs, and player records are increasingly available at granular resolutions. Cricket, and especially the Indian Premier League (IPL), offers a compelling case because the T20 format compresses strategic decisions into a short time window, making batting tempo, bowling discipline, and phase-specific execution central to match outcomes.

Although cricket analytics has received substantial attention in score prediction, winner prediction, and player evaluation studies~\cite{ref1,ref2,ref3,ref4,ref5,ref14,ref15,ref16}, a large portion of applied work remains oriented toward predictive modeling alone. In classroom and practitioner settings, another common limitation is the direct use of raw scorecard data without explicitly addressing duplicates, no-result matches, franchise renaming, super-over contamination, and inconsistent ball-validity logic. Such omissions can distort rate-based statistics, weaken reproducibility, and reduce confidence in downstream inference. Visualization can reveal patterns, but visualization without preprocessing discipline and statistical validation risks overstating informal observations.

This paper addresses that gap by proposing a framework that explicitly unifies three stages: data preprocessing, visual exploration, and hypothesis testing. Instead of presenting IPL analytics as a sequence of isolated charts, we treat preprocessing as a methodological contribution that improves interpretability and inferential reliability. We then combine descriptive graphics with non-parametric tests, confidence intervals, and effect-size aware interpretation.

The main contributions of this study are fourfold:
\begin{enumerate}
    \item A reproducible preprocessing pipeline for IPL 2008--2019 ball-by-ball data, including team normalization, super-over exclusion, legal-delivery filtering, and domain-correct batting/bowling statistics.
    \item A phase-aware visual analysis of batting and bowling behavior across seasons, innings, and match outcomes.
    \item An inferential layer that tests whether key visual patterns are statistically significant rather than merely apparent.
    \item A compact but extensible framework suitable for academic coursework, sports-analytics studies, and future predictive extensions.
\end{enumerate}

\section{Literature Review}
Cricket analytics has evolved from score forecasting and winner prediction toward richer player-evaluation and tactical-analysis pipelines. Early machine-learning work explored score and winning prediction from historical cricket data~\cite{ref1}, while later studies examined T20 outcome prediction and player categorization using supervised learning frameworks~\cite{ref2,ref3,ref4}. Recent IEEE work has extended this direction into broader cricket performance systems and review-oriented analytics pipelines~\cite{ref5,ref6,ref7}. These studies show the growing maturity of cricket analytics, yet they largely emphasize predictive capability over preprocessing transparency and inferential validation.

Beyond cricket, sports analytics literature highlights the value of machine learning, outcome modeling, and historical performance mining across domains~\cite{ref8,ref9,ref10}. Such work motivates structured feature engineering and performance comparison, but cricket remains distinctive because delivery-level sequencing and phase-based tactics strongly influence interpretation. In Twenty20 cricket specifically, player evaluation and tactical aggressiveness have been studied through role-sensitive performance measures~\cite{ref14,ref15}. Network-based cricket analysis has further shown that aggregate runs and wickets alone can miss the contextual quality of matchups~\cite{ref16}.

The visualization component of this paper is informed by foundational work on graphical perception and statistically responsible tabular structure. Cleveland and McGill showed that visual encoding choices materially affect quantitative interpretation~\cite{ref11}, while Heer and Bostock demonstrated empirical evaluation of visual perception tasks in modern data-display settings~\cite{ref12}. Wickham's tidy-data framework provides the structural basis for reliable transformation, grouping, and reshaping of event-level records before visualization or testing~\cite{ref13}. Collectively, these studies support the core premise of this paper: effective sports analytics depends on the joint quality of preprocessing, visual encoding, and inferential reasoning.

Despite these advances, three gaps remain visible in the literature. First, many cricket studies focus on outcome prediction rather than descriptive and inferential insight. Second, preprocessing choices are often underreported, despite their direct impact on statistics such as strike rate, economy, wickets, and matchup matrices. Third, fewer applied studies combine phase-aware visualization with formal hypothesis testing over a cleaned IPL ball-by-ball corpus. This paper positions itself in that gap by examining season-wise scoring evolution, phase-wise scoring variation, win-loss batting and bowling differences, specialist phase behavior, and the analytical reliability gains achieved through preprocessing.

\section{Dataset Description}
The study uses two structured IPL datasets covering seasons 2008--2019:
\begin{itemize}
    \item \texttt{matches.csv}: match-level records with season, venue, toss, result, and winner fields.
    \item \texttt{deliveries.csv}: ball-by-ball records containing batting team, bowling team, batsman, bowler, extras, wickets, and dismissal information.
\end{itemize}

The raw corpus contains 179,078 deliveries with 21 columns and 756 matches with 18 columns. After cleaning, the analysis corpus contains 178,974 deliveries, 173,571 legal deliveries, and 752 matches. The cleaned delivery table has 25 attributes after feature engineering.

\begin{table}[ht]
\centering
\caption{Dataset Summary Before and After Preprocessing}
\resizebox{\columnwidth}{!}{%
\begin{tabular}{lcccc}
\toprule
\textbf{Dataset} & \textbf{Rows (Raw)} & \textbf{Cols (Raw)} & \textbf{Rows (Clean)} & \textbf{Cols (Clean)} \\
\midrule
Deliveries & 179,078 & 21 & 178,974 & 25 \\
Legal deliveries & -- & -- & 173,571 & 25 \\
Matches & 756 & 18 & 752 & 17 \\
Player stats & -- & -- & 152 & 10 \\
Bowler stats & -- & -- & 240 & 7 \\
\bottomrule
\end{tabular}}
\label{tab:dataset}
\end{table}

\section{Data Preprocessing and Feature Engineering}
The preprocessing stage is designed to improve completeness, consistency, and domain correctness. The following operations were applied:
\begin{enumerate}
    \item Removed 23 duplicated delivery rows.
    \item Dropped \texttt{umpire3}, which was largely null.
    \item Removed four no-result matches to ensure valid outcome-linked analysis.
    \item Excluded 81 super-over deliveries to avoid contaminating standard 20-over comparisons.
    \item Standardized franchise names, reducing 15 raw team-name variants to 12 normalized franchises.
    \item Added season information to deliveries by merging match metadata.
    \item Derived a match-phase label: Powerplay (1--6), Middle (7--15), and Death (16--20).
    \item Created helper indicators for dot balls and boundaries.
    \item Constructed a legal-deliveries table by excluding wides, since wides do not consume a legal ball for batsman-facing rate calculations.
\end{enumerate}

Feature engineering then produced two analytical summary tables:
\begin{itemize}
    \item \textbf{player\_stats}: total runs, total balls, fours, sixes, innings, dismissals, strike rate, average, and rule-based batting role.
    \item \textbf{bowler\_stats}: runs conceded, balls, matches, wickets, economy, and overs.
\end{itemize}

This stage improved dataset completeness from 86.36\% to 88.54\%, uniqueness from 99.99\% to 100.00\%, and overall data-quality score from 95.44\% to 96.17\%. Fig.~\ref{fig:data_quality} shows that the largest gain was achieved in completeness while preserving already high validity and uniqueness. Fig.~\ref{fig:team_normalization} demonstrates why this preprocessing step matters analytically: franchise-name consolidation removes categorical fragmentation that would otherwise distort team-level aggregation across seasons.

\optionalfigure{fig01_data_quality.png}
{Comparison of raw and cleaned dataset quality across completeness, validity, uniqueness, and overall score. The figure shows that preprocessing primarily improved completeness while preserving already strong validity and uniqueness, thereby strengthening the reliability of downstream analysis.}
{fig:data_quality}{}

\optionalfigure{fig02_team_normalization.png}
{Franchise normalization before and after cleaning. The figure shows how inconsistent historical team labels are consolidated into a stable franchise vocabulary, improving cross-season comparability and the correctness of team-level aggregation.}
{fig:team_normalization}{}

\section{Methodology}
\subsection{Analytical Design}
The framework combines descriptive analytics, visual analytics, and inferential testing. Preprocessing prepares a trustworthy event log, descriptive summaries reveal broad tendencies, and hypothesis tests evaluate whether observed differences are statistically meaningful. Specifically, the methodology evaluates season-wise scoring evolution, phase-wise scoring variation, win-loss batting and bowling differences, specialist behavior across match phases, and the extent to which preprocessing strengthens downstream analysis.

\subsection{Derived Metrics}
Four core cricket metrics were used:
\begin{equation}
\text{Run Rate} = \frac{\text{Total Runs}}{\text{Legal Balls}} \times 6
\end{equation}

\begin{equation}
\text{Strike Rate} = \frac{\text{Batsman Runs}}{\text{Balls Faced}} \times 100
\end{equation}

\begin{equation}
\text{Economy Rate} = \frac{\text{Runs Conceded}}{\text{Legal Balls Bowled}} \times 6
\end{equation}

\begin{equation}
\text{Boundary \%} = \frac{\text{Boundary Deliveries}}{\text{Legal Balls}} \times 100
\end{equation}

\subsection{Statistical Testing Strategy}
Because several cricket-performance distributions are non-Gaussian and phase-group sizes differ, non-parametric tests were prioritized:
\begin{itemize}
    \item Spearman correlation for season-wise monotonic trends.
    \item Kruskal--Wallis test for multi-phase run-rate comparison.
    \item Mann--Whitney U tests for win/loss comparisons.
    \item Wilcoxon signed-rank test for paired bowler economy comparison between Powerplay and Death phases.
    \item Spearman correlation for association between dot-ball percentage and boundary percentage.
\end{itemize}

\section{Results and Discussion}
\subsection{Season and Innings Trends}
Across the full period, the mean first-innings score was 162.37 runs (95\% CI: 160.16--164.58), while the mean second-innings score was 148.96 runs (95\% CI: 146.69--151.22). On average, teams batting first scored 13.41 more runs than chasing teams. The 2018 and 2019 seasons exhibited the highest batting environments, with 2018 first-innings average reaching 179.8 and 2019 first-innings average reaching 173.6.

\optionalfigure{fig03_season_scoring_trend.png}
{Season-wise average scores for first and second innings from 2008 to 2019. The figure highlights long-term batting inflation, persistent first-innings advantage, and a reduced gap between setting and chasing totals in later seasons.}
{fig:season_trend}{}

Fig.~\ref{fig:season_trend} provides the descriptive basis for the trend analysis: both innings exhibit a clear upward movement over time, with especially strong scoring conditions in the final two seasons. The shaded gap between the two lines further indicates that batting first retained a structural advantage throughout the study period, even though chasing performance improved in later years.

\subsection{Phase-Aware Scoring}
Phase-wise descriptive summaries show that Death overs are the most explosive scoring period. Median run rates were 7.50 in the Powerplay, 7.56 in the Middle overs, and 9.60 in the Death overs. Dot-ball pressure was highest in the Powerplay (47.5\%) and lowest in the Death overs (28.7\%), while boundary frequency was highest in the Death overs (19.5\%).

\optionalfigure{fig04_phase_run_rate.png}
{Run rate by match phase. The figure shows that the principal scoring escalation occurs in Death overs, whereas Powerplay and Middle-over scoring remain comparatively similar, indicating that late-innings acceleration drives most of the phase-based scoring separation.}
{fig:phase_rr}{}

As shown in Fig.~\ref{fig:phase_rr}, phase dynamics are asymmetric rather than gradually increasing. The most substantial tactical shift occurs in overs 16--20, where teams convert accumulated innings stability into aggressive scoring. The closeness of Powerplay and Middle-over run rates is also important, because it suggests that the strongest structural distinction in IPL innings is not between early and middle overs, but between those overs and the Death phase.

\optionalfigure{fig07_batting_phase_heatmap.png}
{Phase-wise strike-rate heatmap for the top batsmen in the study. Higher intensity cells indicate stronger scoring acceleration, especially in the Death overs, and help distinguish specialist finishers from more phase-balanced batters. For example, AB de Villiers (223.8) and CH Gayle (214.5) exhibit exceptional Death-over strike rates, while players such as G Gambhir remain comparatively moderate across phases.}
{fig:batting_phase_heatmap}{}

Fig.~\ref{fig:batting_phase_heatmap} adds player-level nuance to the phase analysis. Players such as AB de Villiers and CH Gayle exhibit exceptionally high Death-over strike rates, confirming that aggregate batting performance alone is insufficient to capture phase specialization. By contrast, batters such as G Gambhir and S Dhawan remain comparatively moderate across phases, indicating a more control-oriented scoring profile. The heatmap therefore strengthens the argument that phase-aware evaluation is more informative than overall strike rate alone.

More specifically, the heatmap shows that AB de Villiers reaches a Death-over strike rate of 223.8 and CH Gayle reaches 214.5, values far above their corresponding Powerplay rates. MS Dhoni also shows a pronounced transition from a conservative Powerplay rate of 80.0 to a Death-over rate of 178.8, which is characteristic of a finisher who maximizes impact late in the innings. In contrast, players such as G Gambhir display relatively flatter progression across phases, reinforcing the distinction between acceleration specialists and stability-oriented top-order batters.

\subsection{Leading Batting and Bowling Performers}
The highest aggregate run-scorers in the cleaned player table were V Kohli (5,421), SK Raina (5,382), RG Sharma (4,902), and DA Warner (4,722). Top wicket-takers were SL Malinga (170), A Mishra (156), PP Chawla (149), and DJ Bravo (147). The matchup matrix between top batsmen and top bowlers further showed that some pairings were highly asymmetric, such as SK Raina against PP Chawla (175 runs) and V Kohli against A Mishra (158 runs), supporting the value of matchup-aware analytics over raw career aggregates.

\optionalfigure{fig08_matchup_heatmap.png}
{Matchup heatmap of runs scored by top batsmen against top bowlers. Darker cells indicate pairings where the batsman accumulated substantially more runs, while lighter cells indicate relatively stronger control from the bowler, revealing interaction effects hidden by aggregate statistics.}
{fig:matchup_heatmap}{}

Fig.~\ref{fig:matchup_heatmap} adds an interaction-level perspective to the study. Rather than relying solely on career aggregates, it reveals that player performance is partly matchup-dependent, which is especially relevant in tactical T20 decision-making. For example, the figure shows strong accumulations for pairings such as SK Raina against PP Chawla and V Kohli against A Mishra, while other cells remain relatively light, indicating more restrictive bowler control. This figure is particularly useful because it moves the analysis beyond simple leaderboards toward matchup-sensitive strategy.

\subsection{Hypothesis Testing Results}
Table~\ref{tab:tests} summarizes the key inferential findings.

\begin{table}[ht]
\centering
\caption{Core Hypothesis-Testing Results}
\resizebox{\columnwidth}{!}{%
\begin{tabular}{p{2.2cm}p{1.9cm}p{2.7cm}p{2.2cm}}
\toprule
\textbf{Question} & \textbf{Test} & \textbf{Result} & \textbf{Interpretation} \\
\midrule
Season trend in first-innings scoring & Spearman & $\rho=0.748$, $p=0.005$ & Significant upward scoring trend \\
Season trend in second-innings scoring & Spearman & $\rho=0.748$, $p=0.005$ & Significant upward chasing trend \\
Run-rate difference across phases & Kruskal--Wallis & $H=586.64$, $p<0.001$, $\epsilon^2=0.132$ & Strong phase effect \\
Powerplay vs Middle run rate & Mann--Whitney U & $p=0.802$ & No meaningful difference \\
Powerplay vs Death run rate & Mann--Whitney U & $p<0.001$, $|r_{rb}|=0.439$ & Death overs significantly higher \\
Middle vs Death run rate & Mann--Whitney U & $p<0.001$, $|r_{rb}|=0.460$ & Death overs significantly higher \\
Batting strike rate in wins vs losses & Mann--Whitney U & $p<0.001$, $|r_{rb}|=0.429$ & Winning teams bat faster \\
Bowling economy in wins vs losses & Mann--Whitney U & $p<0.001$, $|r_{rb}|=0.433$ & Winning teams bowl more economically \\
Bowling wickets in wins vs losses & Mann--Whitney U & $p<0.001$, $|r_{rb}|=0.613$ & Winning teams take more wickets \\
Powerplay vs Death bowler economy & Wilcoxon & $p<0.001$ & Death overs are harder to bowl \\
Dot-ball vs boundary association & Spearman & $\rho=-0.265$, $p<0.001$ & Pressure and boundary rate are inversely related \\
\bottomrule
\end{tabular}}
\label{tab:tests}
\end{table}

\subsection{Season-Wise Scoring Trend}
The positive and statistically significant Spearman correlations for both innings indicate that IPL scoring environments became progressively more batter-friendly over time. Linear trend estimates further suggest an increase of roughly 1.74 runs per season in first innings and 1.84 runs per season in second innings. This supports the visual observation that modern IPL teams are not only setting bigger totals but also chasing more effectively.

\subsection{Phase-Wise Differences}
The Kruskal--Wallis result confirms that phase-level scoring is not homogeneous. However, the pairwise tests reveal a more nuanced picture than a simple ``early overs vs middle overs'' narrative. Powerplay and Middle-over scoring are statistically indistinguishable ($p=0.802$), whereas both differ sharply from the Death overs ($p<0.001$ in each comparison). This implies that the primary structural shift in T20 scoring occurs in overs 16--20.

\optionalfigure{fig10_phase_economy_boxplot.png}
{Distribution of bowler economy by match phase. Death overs show both the highest median economy and the widest spread, indicating that they are not only more expensive on average but also less predictable and more volatile. The median rises from 7.58 in the Powerplay and 7.67 in the Middle overs to 9.49 in the Death overs, with several high-economy outliers above 13 runs per over.}
{fig:economy_box}{}

The bowling-side evidence in Fig.~\ref{fig:economy_box} complements the batting-side phase result. Higher spread and higher central tendency in Death-over economy confirm that this phase concentrates both opportunity and risk. The boxplot also shows that Powerplay and Middle-over distributions overlap substantially, whereas the Death-over distribution shifts upward, reinforcing the inferential result that the main structural break occurs late in the innings. In other words, bowlers do not merely become slightly less economical in the Death overs; the entire performance distribution worsens and becomes more unstable.

The medians printed inside the boxplots make this contrast explicit: Powerplay and Middle overs have nearly identical median economy values of 7.58 and 7.67, while the Death-over median rises sharply to 9.49. The Death-over box is also visibly taller and accompanied by several high-economy outliers above 13 runs per over, indicating that not only is the expected cost higher in this phase, but the risk of extreme punishment is materially greater as well.

\subsection{Winning vs Losing Performance}
Winning team innings had a mean batting strike rate of 137.64 (95\% CI: 135.95--139.32), compared with 119.65 (95\% CI: 117.93--121.36) for losing team innings. On the bowling side, winning teams posted a mean economy of 7.36 (95\% CI: 7.25--7.46), whereas losing teams conceded at 8.45 (95\% CI: 8.35--8.55). Winning bowling units also took substantially more wickets on average (7.13 vs. 4.50). Taken together, these results indicate that IPL victories are associated with both acceleration in batting and disruption in bowling.

\optionalfigure{fig12_bowler_match_impact.png}
{Bowler match-impact comparison using wickets taken in winning matches against total wickets. The figure identifies bowlers whose wicket-taking is more tightly aligned with team success than raw career totals alone would suggest, thereby offering a more outcome-sensitive view of bowling impact.}
{fig:bowler_impact}{}

Fig.~\ref{fig:bowler_impact} is especially valuable for interpretation because it distinguishes high-volume wicket-takers from bowlers whose wickets occur more often in successful match contexts. Bowlers such as SL Malinga, Harbhajan Singh, RA Jadeja, and MM Sharma show a large share of wickets in victories, which suggests that their breakthroughs were not merely frequent but also strongly associated with match-winning performances. This makes the figure more decision-relevant than a simple wicket leaderboard and supports the broader claim that bowling impact should be judged in context, not only by total wickets.

\subsection{Specialist Bowling Behavior}
The paired Powerplay-vs-Death economy test on bowlers with observations in both phases showed a strong shift in difficulty, with median economy rising from 7.80 in the Powerplay to 9.84 in the Death overs. Descriptive specialist tables reinforce this conclusion: bowlers such as R Ashwin, SP Narine, and B Kumar stand out in the Powerplay, while DE Bollinger, R Ashwin, SP Narine, and SL Malinga emerge as elite Death-over performers. Notably, SP Narine appears in both specialist lists, indicating uncommon cross-phase adaptability.

\subsection{Interpretation of Findings}
The results support three central arguments. First, preprocessing materially affects the validity of cricket analytics. Without removing super overs, normalizing team names, and filtering legal deliveries, comparisons involving strike rate, economy, franchise-level summaries, and win-linked analysis would be noisier and potentially misleading. The data-quality gains observed in preprocessing are therefore methodologically important rather than cosmetic.

Second, the inferential results refine the visual story. A visualization-only reading may suggest that scoring differs gradually across all phases, but the hypothesis tests show that the real break occurs between the Death overs and the earlier two phases. This distinction matters for tactical framing, because it implies that successful T20 batting is not uniformly aggressive throughout an innings; instead, teams appear to preserve a distinct acceleration zone late in the innings.

Third, match outcomes are jointly shaped by batting tempo and bowling disruption. The win/loss comparisons reveal that successful teams both score faster and suppress opponents more effectively. This is consistent with T20's short-horizon structure, where even moderate shifts in boundary rate, wicket-taking, and death-over control can strongly alter the result.

\section{Limitations and Future Scope}
This study has four main limitations. First, it covers IPL seasons only up to 2019. Second, the analysis does not integrate venue, toss-context interaction, weather, or pitch-condition covariates. Third, player roles are derived from rule-based thresholds rather than learned clustering models. Fourth, the framework is descriptive-plus-inferential and does not yet incorporate causal modeling or out-of-sample predictive validation as the main objective.

Future work can extend the framework in several directions:
\begin{itemize}
    \item Merge later IPL seasons and related tournament data.
    \item Add venue, opposition, and toss-adjusted mixed-effects models.
    \item Incorporate unsupervised role discovery for batters and bowlers.
    \item Extend matchup analytics using richer contextual features such as batting handedness and bowling style.
    \item Build predictive layers for win probability, expected score, and player-impact forecasting on top of the cleaned analytical base.
\end{itemize}

\section{Conclusion}
This paper proposed a structured framework for IPL analytics that integrates preprocessing, visualization, and hypothesis testing rather than treating them as disconnected steps. Using cleaned IPL ball-by-ball data from 2008--2019, the study showed that scoring has increased significantly across seasons, Death overs are statistically distinct from earlier phases, and winning outcomes are strongly associated with superior batting strike rate, lower bowling economy, and higher wicket counts. Just as importantly, the paper demonstrated that careful preprocessing is a prerequisite for trustworthy cricket analytics. The resulting framework is suitable both for academic study and for practical expansion into predictive sports-intelligence systems.

\section*{Acknowledgment}
The authors express their gratitude to the Department of Artificial Intelligence and Machine Learning, Vishwakarma Institute of Technology, Pune, for the academic environment and support provided for this work.

\begin{thebibliography}{00}

\bibitem{ref1}
T. Singh, V. Singla, and P. Bhatia, ``Score and winning prediction in cricket through data mining,'' in \textit{Proc. Int. Conf. Soft Comput. Techniques Implementations}, 2015. [Online]. Available: \url{https://ieeexplore.ieee.org/abstract/document/7489605}

\bibitem{ref2}
``Analysis and Winning Prediction in T20 Cricket using Machine Learning,'' \textit{IEEE Conf. Publication}, 2022. [Online]. Available: \url{https://ieeexplore.ieee.org/document/9807929}

\bibitem{ref3}
``Cricket Players Performance Prediction and Evaluation Using Machine Learning Algorithms,'' \textit{IEEE Conf. Publication}, 2023. [Online]. Available: \url{https://ieeexplore.ieee.org/document/10127503/}

\bibitem{ref4}
R. Suguna, Y. Praveen Kumar, J. Suriya Prakash, P. S. Neethu, and S. Kiran, ``Utilizing Machine Learning for Sport Data Analytics in Cricket: Score Prediction and Player Categorization,'' in \textit{Proc. 2023 IEEE 3rd Mysore Sub Section Int. Conf.}, 2023. [Online]. Available: \url{https://ieeexplore.ieee.org/document/10396955}

\bibitem{ref5}
``Cricket Performance Analysis System Using Machine Learning Techniques,'' \textit{IEEE Conf. Publication}, 2024. [Online]. Available: \url{https://ieeexplore.ieee.org/document/10864240/}

\bibitem{ref6}
``Cricket Analytics: ODI Match Score Prediction and Resource Metric Modeling using Machine Learning Model,'' \textit{IEEE Conf. Publication}, 2024. [Online]. Available: \url{https://ieeexplore.ieee.org/document/10796074}

\bibitem{ref7}
``World Cup Cricket Data Analytics: A Comparative Review of T20 and ODI World Cup Prediction Techniques,'' \textit{IEEE Conf. Publication}, 2025. [Online]. Available: \url{https://ieeexplore.ieee.org/abstract/document/11034937}

\bibitem{ref8}
``Machine Learning in Sports Analytics: Volleyball Match Outcome Prediction,'' \textit{IEEE Conf. Publication}, 2024. [Online]. Available: \url{https://ieeexplore.ieee.org/document/10796894}

\bibitem{ref9}
``Forecasting Sports Performance Using Historic Data and Machine Learning,'' \textit{IEEE Conf. Publication}, 2024. [Online]. Available: \url{https://ieeexplore.ieee.org/document/10860038}

\bibitem{ref10}
``Relative Analysis and Performance of Machine Learning Approaches in Sports,'' \textit{IEEE Conf. Publication}, 2021. [Online]. Available: \url{https://ieeexplore.ieee.org/document/9676147}

\bibitem{ref11}
W. S. Cleveland and R. McGill, ``Graphical perception: Theory, experimentation, and application to the development of graphical methods,'' \textit{J. Amer. Statist. Assoc.}, vol. 79, no. 387, pp. 531--554, 1984.

\bibitem{ref12}
J. Heer and M. Bostock, ``Crowdsourcing graphical perception: Using Mechanical Turk to assess visualization design,'' in \textit{Proc. SIGCHI Conf. Human Factors in Computing Systems}, 2010, pp. 203--212.

\bibitem{ref13}
H. Wickham, ``Tidy data,'' \textit{J. Stat. Softw.}, vol. 59, no. 10, pp. 1--23, 2014.

\bibitem{ref14}
C. J. Petersen, ``Analysis of Twenty/20 cricket performance during the 2008 Indian Premier League,'' \textit{Int. J. Perform. Anal. Sport}, vol. 8, no. 3, pp. 63--69, 2008. [Online]. Available: \url{https://www.tandfonline.com/doi/abs/10.1080/24748668.2008.11868448}

\bibitem{ref15}
J. Davis, H. Perera, and T. B. Swartz, ``Player evaluation in Twenty20 cricket,'' \textit{Aust. N. Z. J. Stat.}, vol. 57, no. 1, pp. 55--71, 2015. [Online]. Available: \url{https://journals.sagepub.com/doi/10.3233/JSA-150002}

\bibitem{ref16}
S. Mukherjee, ``Quantifying individual performance in cricket -- A network analysis of batsmen and bowlers,'' \textit{Physica A}, vol. 393, pp. 624--637, 2014. [Online]. Available: \url{https://www.sciencedirect.com/science/article/pii/S0378437113008819}

\bibitem{ref17}
``The Determinants of Success in One Day International (ODI) and Twenty20 (T20) Cricket Matches: A Systematic Review and Meta-Analysis,'' \textit{Applied Sciences}, 2025. [Online]. Available: \url{https://www.mdpi.com/2076-3417/15/19/10341}

\end{thebibliography}

\end{document}
